\documentclass{resume} % Use the custom resume.cls style

\usepackage[left=0.4 in,top=0.4in,right=0.4 in,bottom=0.4in]{geometry} % Document margins
\usepackage{hyperref}
\usepackage{tcolorbox}
\usepackage{xcolor}
\newenvironment{problem}[2][Problem]{\begin{trivlist}
		\item[\hskip \labelsep {\bfseries #1}\hskip \labelsep {\bfseries #2.}]}{\end{trivlist}}

\newcommand{\tab}[1]{\hspace{.2667\textwidth}\rlap{#1}} 
\newcommand{\itab}[1]{\hspace{0em}\rlap{#1}}
\name{Ching-Lin Huang} % Your name
% You can merge both of these into a single line, if you do not have a website.
\address{huan2919@umn.edu \\ (763)-485-2249 \\ \href{https://scholar.google.com.tw/citations?user=mytujjwAAAAJ&hl=zh-TW}{google scholar} \\ \href{https://github.com/ChingLinHuang}{github}} 


\begin{document}
	
	
	\begin{rSection}{Education}
		
		\begin{minipage}[t]{0.8\textwidth}
			\textbf{Ph.D., Ecology, Evolution, and Behavior, University of Minnesota, Twin Cities}\\
			Advisors: Dr. Allison K Shaw
		\end{minipage}
		\begin{minipage}[t]{0.2\textwidth}
			\hfill {2024-}
		\end{minipage}
		
		\begin{minipage}[t]{0.8\textwidth}
			\textbf{M.S., Institute of Ecology and Evolutionary Biology, National Taiwan University, Taipei, Taiwan}\\
			Thesis title: \textit{"Disentangling ecological processes underlying the metacommunity by integrating several analytical methods"}; Advisors: Dr. David Zelený
		\end{minipage}
		\begin{minipage}[t]{0.2\textwidth}
			\hfill {2020-2023}
		\end{minipage}
		
		\begin{minipage}[t]{0.8\textwidth}
			\textbf{B.S., Department of Mathematics, National Taiwan University, Taipei, Taiwan}
		\end{minipage}
		\begin{minipage}[t]{0.2\textwidth}
			\hfill {2016-2020}
		\end{minipage}
		
		
		
	\end{rSection}
	
	
	\begin{rSection}{PUBLICATIONS AND CONFERENCES}
		\textbf{Manuscripts}
		\begin{itemize}
			\item \underline{\textbf{Huang, C.-L.}}, Zelený, D., Chang-Yang, C.-H. (2023) ``Integrating several analytical methods to assess strength of ecological processes behind metacommunity assembly." \textit{Oikos}, e10166, \url{https://doi.org/10.1111/oik.10166}
			\item \underline{\textbf{Huang, C.-L.}}, Wan, J., Ke, P.-J., Wei, S., Chang-Yang, C.-H. (Accepted in \textit{Plant and Soil}) ``Plant–soil feedback persists beyond host death to shape density-dependent plant competition."
		\end{itemize}

	
		\textbf{Conferences}
		\begin{itemize}
			\item \underline{\textbf{Huang, C.-L.}}, Zelený, D. (2022) ``How well do available methods quantify the importance of assembly processes in ecological community data: a simulation approach." {\color{gray}The 10th Biennial conference of the International Biogeography Society}
			\item \underline{\textbf{Huang, C.-L.}}, Shaw, A. (2025) ``Frugivore-plant interactions affect plant population persistence and spreading speed, mediated by an Allee effect". {\color{gray} Ecological Society of American 2025 Annual Meeting}
			\item \underline{\textbf{Huang, C.-L.}}, Shaw, A., Wright, A. (2025) ``Species-driven microhabitat modifications influence coexistence and productivity under resource competition". {\color{gray} Society for Modeling \& Theory in Population Biology 2026 Meeting}
			
			
		\end{itemize}
	\end{rSection}


		\begin{rSection}{AWARDS}
		\begin{itemize}
			\item Honorable Mention of Lotka Award in ESA 2025
		\end{itemize}
	\end{rSection}
	
	\begin{rSection}{WORKING EXPERIENCE}
		\textbf{Research assistant}\\
		\textbf{- Stanton Lab, University of Minnesota, Twin Cities} \ PI: Dr. Daniel Stanton \hfill{Fall 2025}\\
		\textbf{- Ke Lab, National Taiwan University} \ PI: Dr. Po-Ju Ke \hfill{2023-2024}
%		\begin{tcolorbox}
%			\small
%			\begin{itemize}
%				\item Studying how the interactive effects of soil microbial conditioning and decay time determine the competitive outcome and stability of plant community
%				\item Greenhouse experiments to explore how and how long the microbial conditioning after plants die would influence plant growth
%				\item Theoretical modeling to explore how age/stage-dependent plant-soil microbe interactions would affect competition consequences
%			\end{itemize}
%		\end{tcolorbox}
		

		
		\textbf{Part-time research assistant}\\
		\textbf{- Vegetation Ecology Lab, National Taiwan University} \ PI: Dr. David Zelený \hfill{2020-2022}
%		\begin{tcolorbox}
%			\small
%			\begin{itemize}
%				\item Studying how fog frequency affects plant taxonomic and functional turnover in multiple scales
%				\item Field data collection, including overstory and understory census and leaf traits collection
%				\item Identifying fog specialist species by analyzing data from herbarium and vegetation survey and remote sensing
%			\end{itemize}
%		\end{tcolorbox}

		
		\textbf{Teaching assistant}\\
		\textbf{- Ecology: Theory and Concept} \ Lecturer: Dr. David Tilman \hfill{Fall 2025}\\
		\textbf{- Introduction to Theoretical Ecology (\href{https://chinglinhuang.github.io/2022_Fall_TheoEcol/}{Course Link})} \ Lecturer: Dr. Po-Ju Ke \hfill{Fall 2022}\\ 
		\textbf{- Numerial Methods in Community Ecology} \ Lecturer: Dr. David Zelený \hfill{Spring 2022}\\
		\textbf{- General Botany Lab} \ Lecturer: Dr. Fen-Ming Lee \hfill{Fall 2021}\\
		\textbf{- Introduction to R for Ecologists} \ Lecturer: Dr. David Zelený \hfill{Fall 2020}\\
		\textbf{- Calculus} \ Lecturer: Dr. Ya-Ju Tsai \hfill{Fall 2019- Spring 2020}\\
		
		
		
	\end{rSection}
	
	\newpage
	\begin{rSection}{FUNDING}
		EEB Career Development Funding (\$2,000) \hfill{Spring 2025}\\
		EEB Research Funding (\$1,500) \hfill{Summer 2025, 2026}\\ 
		EEB Travel Funding (\$900) \hfill{Summer 2025, 2026}\\
		SMTPB Conference Lodging Support \hfill{Summer 2026}
	\end{rSection}

	
	\begin{rSection}{OTHER EXPERIENCE}
		\textbf{Journal manuscript referee}\\
		\textit{- Journal of Vegetation Science} (1 review)
		
		\textbf{Selected field survey and greenhouse experience}\\
		- Bio3D greenhouse experiment \hfill{2025}\\
		- Plant--soil feedback density gradient experiment \hfill{2023--2024}\\
		- Overstory and understory composition survey and leaf traits collection in fog frequency \\ \hspace*{0.5em} gradient transect in Northern Taiwan \hfill{2021-2022}\\
		- Leaf trait collection, and seedling survey in Fushan Forest Dynamics Plot \hfill{2019-2020}\\
		- Tree census in Lalashan Forest Dynamics Plot \hfill{2019-2020}\\
		- Tree and seedling census in Kenting Forest Recreation Area \hfill{Summer 2019}\\
		- Tree census in Fushan Forest Dynamics Plot \hfill{Winter 2019}
		
		
	\end{rSection}
	
\end{document}
